\section{Joint-metric selection on locked external tasks}
\label{sec:multidimensional-rerun}

We extended the label-free trajectory record with a 24-dimensional behavioral profile and a phase $\times$ action-family $\times$ channel tensor of shape $4\times8\times3$. The profile measures command and action diversity, repetition, local transitions, run structure, and within-trace distribution shift. The tensor preserves the phase, action-family, and channel axes; training-fold HOSVD provides a lower-rank comparison. We also evaluated IDF-weighted command and ordered-bigram overlap. The link score uses observed behavior only: Human/Agent labels are stored separately and used for evaluation, never as pair features. A prior development error-case catalog showed 208 introduced cross-task Human--Agent links versus 7 removed links when predicted execution-mode similarity was added; this motivated excluding mode similarity from the association score.

Before scoring, a schema audit found that the selected CyBashBench and NL2Bash Claude Sonnet sets contain prompt/final-answer records without bash execution events. We therefore exclude them from the primary comparison rather than treating generated answers as executed trajectories. The locked external primary set contains 10 CVEBench and 4 CyberGym benchmark-qualified task groups from Cyber Task Horizons, revision \texttt{178c0da1800f9c07311e1b8d6c61dd8104e7b4ae} (\href{https://github.com/lyptus-research/cyber-task-horizons-data/tree/178c0da1800f9c07311e1b8d6c61dd8104e7b4ae}{CC BY 4.0 source}). Thus the external inference covers two benchmarks and 14 tasks. Same-task labels are retrieval proxies, not actor or organization identity labels.

Development selection used task-grouped nested folds and a joint direction-adjusted objective over the observable error, coverage, and ranking metrics. The final development-set selector chose IDF partial sequence, while all three nested outer-fold selectors chose command TF--IDF, indicating unstable method selection. The locked external set covers 10 CVEBench and 4 CyberGym task groups; Agent--Agent false-split rate was not estimable because no task had multiple Agent trajectories.

\begin{table}[t]
\centering
\small
\setlength{\tabcolsep}{6pt}
\begin{tabularx}{\linewidth}{l r X r}
\toprule
Metric & Baseline & IDF partial sequence & Paired $\Delta$ (95\% CI), pp \\
\midrule
False-merge rate & 100.0 & 95.3 & $-4.7\;[-7.2,-2.2]$ \\
False-split rate & 0.0 & 0.0 & $0.0\;[0.0,0.0]$ \\
Human--Agent false merges & 100.0 & 97.2 & $-2.8\;[-4.9,-0.8]$ \\
Human--Agent false splits & 0.0 & 0.0 & $0.0\;[0.0,0.0]$ \\
Human--Human false merges & 100.0 & 87.3 & $-12.7\;[-21.5,-3.9]$ \\
Human--Human false splits & 0.0 & 0.0 & $0.0\;[0.0,0.0]$ \\
Agent--Agent false merges & 100.0 & 100.0 & $0.0\;[0.0,0.0]$ \\
Agent--Agent false splits & -- & -- & Not estimable \\
Balanced accuracy & 50.0 & 52.3 & $+2.3\;[+1.1,+3.6]$ \\
Query coverage & 100.0 & 100.0 & $0.0\;[0.0,0.0]$ \\
Hit@1 & 0.0 & 10.4 & $+10.4\;[0.0,+23.3]$ \\
MRR & 18.3 & 27.9 & $+9.6\;[+1.1,+19.0]$ \\
Recall@5 & 40.4 & 50.0 & $+9.6\;[-12.5,+29.2]$ \\
NDCG@5 & 18.7 & 28.7 & $+10.1\;[-3.5,+19.4]$ \\
\bottomrule
\end{tabularx}
\caption{External task-macro rates and ranking scores are percentages. Paired changes are method minus command TF--IDF; negative changes improve error rates, while positive changes improve accuracy, coverage, and ranking. Confidence intervals resample paired benchmark-qualified tasks within benchmark (5,000 draws).}
\label{tab:multidimensional-rerun}
\end{table}

The zero FSR is real for this frozen operating point, but it does not indicate useful discrimination. The external set contains 36 same-task positive pairs and 482 different-task negative pairs: the selected threshold accepted all 36 positives and 458 negatives (95.0\% micro, 95.3\% task-macro FMR). The command baseline accepted all 518 pairs. Threshold selection constrained development FSR to 20\% but imposed no hard FMR ceiling, so the external result is close to linking every candidate. Balanced accuracy is 52.3\%, near chance. FSR is the fraction of positive pairs rejected; FMR is the fraction of negative pairs accepted. Both are averaged over queries and tasks. The relative FMR and joint-utility criteria pass, but the absolute error rate makes this operating point unusable for attribution.

\begin{figure}[tb]
  \centering
  \includegraphics[width=\linewidth]{multidimensional-rerun-v1.pdf}
  \caption{Direction-adjusted external changes for five candidate methods relative to command TF--IDF. Positive values improve the metric; error rates are sign-reversed. The selected IDF partial-sequence row is bold. An asterisk marks a paired 95\% task-cluster interval supporting improvement.}
  \label{fig:multidimensional-rerun}
\end{figure}

The per-trace 24-dimensional profiles and $4\times8\times3$ tensors were retained as separate feature records, with evaluation labels in a separate file. Neither profile distance nor tensor projection was selected by the joint objective; local command and ordered-bigram overlap was selected. This does not show that a tensor profile is an identity signature. Human terminal inputs and Agent tool calls use different capture interfaces, and the matched tasks do not establish common operators, attackers, or organizations.

A separate exploratory check tested whether the 24 scalar features add Human/Agent source-label signal on the same 177 traces. Command TF--IDF reached 84.9\% balanced accuracy; the numeric profile alone reached 65.6\%, and inner-selected profile fusion reached 83.9\%. Its paired change was $-1.0$ percentage point (95\% task-cluster CI $[-3.1,+1.0]$); all six metric intervals included zero. The profile received nonzero weight in 5 of 18 outer folds. This reused cohort provides no evidence of a reliable incremental classification gain.

\begin{figure}[tb]
  \centering
  \includegraphics[width=0.9\linewidth]{behavior-profile-v2.pdf}
  \caption{Direction-adjusted change in six Human/Agent source-label metrics for inner-selected profile fusion versus command TF--IDF. Positive values favor profile fusion; Brier change is sign-reversed so positive indicates lower error. Error bars are paired, benchmark-stratified task-cluster 95\% intervals. All intervals include zero. The 40 task groups had been inspected previously, so this is exploratory reuse.}
  \label{fig:behavior-profile-v2}
\end{figure}
